\documentclass[tikz,border=8pt]{standalone}
\input{../../../docs/tikz/dag_preamble.tex}

\begin{document}

\begin{tikzpicture}[
  x=1cm,
  y=1cm,
  every node/.append style={outer sep=3pt,font=\sffamily\scriptsize},
  every path/.append style={font=\sffamily\scriptsize}
]

\node[align=left,text width=30cm,font=\sffamily\small,anchor=north west]
  at (0,0) {\textbf{Paper:} EFX for Additive Chores\\
  \textbf{Authors:} Wentao He and Biaoshuai Tao\\
  \textbf{Source:} arXiv v2 (2026-07-09)\\
  \textbf{Scope:} named theoretical results and displayed definitions; every
  node below has a compiled paper-facing Lean route.};

\node[dag_model] (EF) at (0,-4) {\textbf{Definition}\\Envy-free for chores};
\node[dag_model] (EFX) at (6,-4) {\textbf{Definition}\\EFX for chores};
\node[dag_model] (PO) at (12,-4) {\textbf{Definition}\\Pareto optimality};
\node[dag_model] (CanonicalDef) at (18,-4) {\textbf{Definition}\\Canonical M01 allocation};
\node[dag_model] (Model) at (24,-4) {\textbf{Additive chore model}\\Finite allocation, additive cost, M01/M2/M34 partition};

\node[dag_lemma] (TriNoTwo) at (0,-8) {\textbf{Tri proposition}\\No bundle has two A items};
\node[dag_lemma] (TriNoA) at (6,-8) {\textbf{Tri proposition}\\The no-A bundle costs at least 20};
\node[dag_lemma] (AppendixNoTwo) at (12,-8) {\textbf{Appendix A}\\No bundle has two A items};
\node[dag_lemma] (AppendixBounds) at (18,-8) {\textbf{Appendix A}\\$p_1,p_2\ge r$};
\node[dag_result] (TriTheorem) at (24,-8) {\textbf{Theorem 1}\\Tri-valued EFX nonexistence for every $n\ge4$};

\node[dag_lemma] (POLarge) at (0,-12) {\textbf{Theorem 2 proposition}\\Every EFX agent has a large chore};
\node[dag_lemma] (POBounds) at (6,-12) {\textbf{Theorem 2 proposition}\\All own costs $\ge r+1$; one $\ge r+2$};
\node[dag_result] (POTheorem) at (12,-12) {\textbf{Theorem 2}\\EFX and Pareto optimality are incompatible};
\node[dag_lemma] (M34) at (18,-12) {\textbf{Lemma}\\M34 insertion preserves EFX};
\node[dag_lemma] (Composition) at (24,-12) {\textbf{Lemma}\\Cost-gap composition gives EFX};

\node[dag_lemma] (CanonicalProps) at (0,-16) {\textbf{Lemma}\\Canonical and super-canonical M01 properties};
\node[dag_lemma] (Balanced) at (6,-16) {\textbf{Lemma}\\Balanced multigraph orientation};
\node[dag_lemma] (M2) at (12,-16) {\textbf{Lemma}\\Every M2 pool has an EFX allocation};
\node[dag_lemma] (M2Props) at (18,-16) {\textbf{Lemma}\\M2 residue properties and preferred endpoints};
\node[dag_lemma] (Exceptional) at (24,-16) {\textbf{Lemma}\\Exceptional-residue combination};

\node[dag_result] (Four) at (12,-20) {\textbf{Theorem 3}\\Every four-agent bi-valued chore instance has a feasible EFX allocation};
\node[dag_lemma] (Low) at (0,-20) {\textbf{Source bridge}\\Low-ratio reverse-first round robin};
\node[dag_lemma] (High) at (6,-20) {\textbf{Source bridge}\\High-ratio M01/M2 modulo dispatch};
\node[dag_lemma] (Normalize) at (18,-20) {\textbf{Source bridge}\\Positive normalization; constant and binary cases};
\node[dag_model] (Library) at (24,-20) {\textbf{Definitions provenance}\\GameTheoryInLean EconCSLib allocation and owner-map design};

\draw[dag_arrow] (EFX) -- (TriNoTwo);

\draw[dag_arrow] (EFX) -- (POLarge);
\draw[dag_arrow] (POLarge) -- (POBounds);
\draw[dag_arrow] (POBounds) -- (POTheorem);
\draw[dag_arrow] (PO) -- (POTheorem);

\draw[dag_arrow] (Balanced) -- (M2);
\draw[dag_arrow] (M2) -- (M2Props);
\draw[dag_arrow] (M2Props) -- (Exceptional);
\draw[dag_arrow] (Composition) -- (Exceptional);

\draw[dag_arrow] (High) -- (Four);
\draw[dag_arrow] (Exceptional) -- (Four);
\draw[dag_arrow] (Normalize) -- (Four);

\node[align=left,text width=30cm,font=\sffamily\scriptsize,anchor=north west]
  at (0,-22.1) {\textit{Reading guide:} nodes are grouped by the three source
  theorem routes; the displayed solid arrows are selected direct Lean-discharged
  dependencies. The ``source bridge'' nodes summarize proved internal dispatches
  rather than external assumptions.};

\end{tikzpicture}
\end{document}
